\definecolor{headergreen}{HTML}{00B08B} % Colore smeraldo dell'intestazione

\begin{center}
\renewcommand{\arraystretch}{1.8}
\resizebox{\linewidth}{!}{%
\begin{tabular}{|c|c|c|c|c|}
    \hline
    \rowcolor{headergreen} 
    \textbf{\textcolor{white}{Messaggio}} & 
    \textbf{\textcolor{white}{Parola di codice}} & 
    \textbf{\textcolor{white}{Singolo errore}} & 
    \textbf{\textcolor{white}{Doppio errore}} & 
    \textbf{\textcolor{white}{Triplo errore}} \\ \hline
    
    % --- Messaggio 00 ---
    \textbf{00} & \textbf{000} & (001 010 100) & (011 110 101) & \textcolor{blue}{111} \\ \hline
    
    % --- Messaggio 01 ---
    \textbf{01} & \textbf{01\textcolor{red}{1}} & (010 001 111) & (000 101 110) & \textcolor{blue}{100} \\ \hline
    
    % --- Messaggio 10 ---
    \textbf{10} & \textbf{10\textcolor{red}{1}} & (100 111 001) & (110 011 000) & \textcolor{blue}{010} \\ \hline
    
    % --- Messaggio 11 ---
    \textbf{11} & \textbf{11\textcolor{red}{0}} & (111 100 010) & (101 000 011) & \textcolor{blue}{001} \\ \hline
\end{tabular}%
}

\vspace{0.5cm}
\begin{minipage}{\linewidth}
    \small
    \textbf{Nota sulle capacità di rilevazione:} \\
    \indent \textbf{Singolo e Triplo errore:} \textbf{Rilevati} perché il numero dispari di errori altera la parità. \\
    \indent \textbf{Doppio errore:} \textcolor{red}{\textbf{Non rilevati}} perché due errori mantengono la parità inalterata (pari).
\end{minipage}
\end{center}